\documentclass[../Main.tex]{subfiles}
\begin{document}

\section{Overall Technology Blocks}
\label{section:3.1}

To address the requirements identified in Chapter~2, the proposed grading
application is organized into four main technology blocks: a frontend
interaction block, a backend and data-management block, a document-processing
and OMR block, and an OCR block for student-name recognition. This organization
is chosen because the application is not only a recognition pipeline, but also a
teacher-oriented web application that must support scanning, review, storage,
and later result inspection in one workflow.

The frontend interaction block is responsible for the teacher-facing interface.
Its purpose is to support class and examination management, answer-sheet image
capture or upload, review of OCR/OMR outputs, submission confirmation, and
later result inspection. The backend and data-management block is responsible
for request handling, workflow coordination, and persistent storage of academic
and grading data.

The document-processing and OMR block is responsible for transforming a raw
answer-sheet image into structured mark-based information. Its purpose is to
align the answer sheet, extract the predefined functional regions, and recognize
student ID, exam code, and selected answers. The OCR block complements this
process by handling the student-name text region, which cannot be processed
reliably by OMR alone. Together, these four blocks form the technological
foundation of the proposed application.

\section{Frontend Interaction Block}
\label{section:3.2}

The frontend interaction block is the part of the application through which the
teacher performs daily grading tasks. Its purpose is to provide an interactive
interface for academic data management, answer-sheet scanning, review of
recognition results, submission handling, and later result inspection. In this
thesis, this block is implemented using React, TypeScript, and Vite.
Accordingly, this block mainly addresses the Chapter~2 requirements for
managing academic data, reviewing recognition results, saving submissions, and
viewing stored grading outcomes.

\subsection{React}
\label{subsection:3.2.1}

React is a JavaScript library for building component-based user interfaces
\cite{react2026}. In this thesis, it belongs to the frontend interaction block,
whose goal is to support teacher-facing operations such as class management,
answer-sheet review, submission inspection, and analytics viewing.

React is used because these views are highly interactive and depend on changing
application state. When the teacher edits an answer key, reviews OCR/OMR
outputs, switches filters, or opens a saved submission, the displayed
information must update immediately and consistently. Within this block, React
provides the component structure that organizes these views and keeps the user
interface aligned with the grading workflow described in Chapter~2.

\subsection{TypeScript}
\label{subsection:3.2.2}

TypeScript is the typed programming language used together with React in the
frontend block~\cite{typescript2026}. In the proposed application, it is mainly
used to represent structured data exchanged between the frontend and the
backend, including classes, students, examinations, exam codes, submissions,
and grading results.

TypeScript is chosen because the frontend does not work with simple unstructured
pages; instead, it manages many related data objects that must remain consistent
across review, result-management, and analytics screens. Within the frontend
interaction block, TypeScript helps define these data structures clearly and
reduces interface-level errors when API payloads are created, received, and
updated.

\subsection{Vite}
\label{subsection:3.2.3}

Vite is the build and development tool used for the frontend interaction block
~\cite{vite2026}. Its role in this thesis is not to provide an end-user feature,
but to support the implementation and maintenance of the React-TypeScript
interface during development.

Vite is selected because the project requires a lightweight and practical
frontend workflow rather than a heavy build environment. Within the frontend
block, it is used to run the development server, bundle the application, and
keep frontend development and testing simple during implementation.

\section{Backend and Data-Management Block}
\label{section:3.3}

The backend and data-management block is responsible for connecting the
teacher-facing interface with the recognition pipeline and the persistent
storage layer. Its purpose is to receive frontend requests, coordinate
processing steps, expose application APIs, and store academic and grading
records. In this thesis, this block is implemented using FastAPI, SQLAlchemy,
and SQLite.
Accordingly, this block mainly addresses the Chapter~2 requirements for
academic-data management, submission storage, and retrieval of saved results
and analytics.

\subsection{FastAPI}
\label{subsection:3.3.1}

FastAPI is a Python web framework for building APIs with structured request and
response handling~\cite{fastapi2026}. In this thesis, it belongs to the backend
and data-management block, whose purpose is to connect the teacher-facing
interface with the recognition and storage layers.

FastAPI is used because the recognition pipeline is also implemented in Python,
so web requests can be connected directly with OCR/OMR processing without
introducing a separate service layer. Within the backend block, FastAPI
provides the endpoints for answer-sheet processing, submission storage,
academic-data management, and result or analytics retrieval. Alternative
backend frameworks such as Express or Spring could also be used, but they would
require additional integration effort because the recognition pipeline is
already implemented in Python.

\subsection{SQLAlchemy}
\label{subsection:3.3.2}

SQLAlchemy is the object-relational mapping tool used in the backend and
data-management block~\cite{sqlalchemy2026}. In the proposed application, it is
used to define and manipulate the main academic entities, including classes,
students, examinations, exam codes, answer keys, and submissions.

SQLAlchemy is chosen because the application needs a structured and maintainable
way to work with persistent data through the backend APIs. Within this block, it
supports model definition, database-session handling, and CRUD operations,
allowing the backend logic to work with structured objects rather than raw SQL
statements.

\subsection{SQLite}
\label{subsection:3.3.3}

SQLite is the persistent storage technology used in the current implementation
~\cite{sqlite2026}. It is part of the backend and data-management block because
this block is responsible for storing academic configuration, saved
submissions, grading results, and related metadata.

SQLite is selected because the proposed application is intended for a single
teacher and does not target large-scale multi-user deployment. Within this
block, it provides a lightweight local database that is easy to set up and
maintain during development and practical use.

\section{Document-Processing and OMR Block}
\label{section:3.4}

The document-processing and OMR block is responsible for handling the
structured visual parts of the answer sheet. Its purpose is to transform a raw
camera-captured or uploaded image into aligned regions from which mark-based
grading information can be extracted. In this thesis, this block is implemented
mainly with OpenCV and rule-based OMR procedures.
Accordingly, this block mainly addresses the Chapter~2 requirements for answer-sheet image
processing and recognition of structured grading fields.

\subsection{OpenCV}
\label{subsection:3.4.1}

OpenCV is a computer-vision library that provides image-processing and
geometric-transformation functions~\cite{bradski2000opencv}. In this thesis, it
belongs to the document-processing and OMR block, whose purpose is to transform
raw answer-sheet images into aligned and analyzable regions.

OpenCV is used because the application must process practical input images that
may contain rotation, perspective distortion, uneven lighting, or background
noise. Within this block, it is used for grayscale conversion, thresholding,
contour detection, perspective transformation, and region-based cropping.
Alternative libraries such as Pillow or scikit-image could support some
preprocessing tasks, but they are less suitable as the main technology because
this application also requires contour-based alignment and perspective
transformation.

\subsection{Rule-Based OMR}
\label{subsection:3.4.2}

Rule-based OMR is the recognition approach used on top of the OpenCV-based
preprocessing pipeline in this block. Its role is to extract structured
mark-based information from the aligned answer sheet, specifically the student
ID, exam code, and selected answers.

Rule-based OMR is chosen because these fields follow predefined layouts in which
each intended value corresponds to a filled mark at a known position. Within
the document-processing and OMR block, the aligned and cropped regions are
divided into cells or answer-choice areas, and the marks are analyzed through
fill evidence and threshold-based decision rules. An end-to-end learning-based
recognizer or specialized OMR hardware could also be considered, but rule-based
OMR is more appropriate here because the supported answer sheet has a fixed
layout and the thesis targets a lightweight application that works with
ordinary captured images.

\section{OCR Block for Student-Name Recognition}
\label{section:3.5}

The OCR block is responsible for the textual part of the answer sheet that
cannot be handled effectively by OMR. Its purpose is to recognize the
student-name region after the answer sheet has already been aligned and cropped
by the document-processing block. In this thesis, this block is implemented
using a YOLO-based text detector, a CRNN-based recognizer, and PyTorch.
Accordingly, this block mainly addresses the Chapter~2 requirement for recognizing textual
student information that does not follow a fixed mark grid.

\subsection{YOLO}
\label{subsection:3.5.1}

YOLO is an object-detection framework designed for efficient visual detection
tasks~\cite{redmon2016yolo}. In this thesis, it belongs to the OCR block, whose
purpose is to recognize the student-name region that cannot be handled
reliably by OMR.

YOLO is used because the student-name region may contain multiple words with
variable spacing and writing position. Within the OCR block, it is used to
localize the word-level text regions inside the cropped student-name area
before recognition is performed. Alternative text-detection approaches such as
CRAFT or Faster~R-CNN could also be considered, but YOLO is chosen because it
offers a practical balance between detection performance and implementation
simplicity for this thesis.

\subsection{CRNN}
\label{subsection:3.5.2}

CRNN is a sequence-recognition model that combines convolutional feature
extraction with recurrent sequence modeling. In the OCR block, it is used to
recognize the text content of the word images detected by YOLO.

CRNN is chosen because OCR on student-name images is not well suited to
explicit character-by-character segmentation, especially under handwriting
variation or image noise. Within this block, the detected word images are
passed to the CRNN recognizer, which predicts the corresponding text sequence.
The theoretical basis of this approach is related to the CTC framework
introduced by Graves et al.~\cite{graves2006ctc} and the CRNN-based
sequence-recognition approach described by Shi et al.~\cite{shi2017crnn}.
General-purpose OCR tools such as Tesseract or cloud OCR services are possible
alternatives, but they are less aligned with the handwritten Vietnamese
name-recognition requirement and the self-contained scope of the application.

\subsection{PyTorch}
\label{subsection:3.5.3}

PyTorch is the deep-learning framework used to implement and run the OCR block
~\cite{paszke2019pytorch}. In this thesis, it provides the runtime environment
for loading and executing the OCR models inside the Python backend.

PyTorch is selected because it supports flexible model implementation and
integrates naturally with the surrounding Python-based backend. Within the OCR
block, it is used to execute the YOLO and CRNN models so that student-name OCR
results can be returned to the main grading workflow together with the OMR
outputs.

\section{Chapter Summary}
\label{section:3.6}

This chapter has presented the technological framework of the proposed grading
application in terms of its main functional blocks. The frontend interaction
block supports teacher-facing operations, the backend and data-management block
coordinates requests and stores persistent records, the document-processing and
OMR block handles structured mark-based regions, and the OCR block handles the
student-name text region.

The technologies used in these blocks are chosen according to their roles in the
application. React, TypeScript, and Vite support the frontend interface;
FastAPI, SQLAlchemy, and SQLite support backend coordination and persistent
storage; OpenCV and rule-based OMR support answer-sheet alignment and
mark-based recognition; and YOLO, CRNN, and PyTorch support student-name OCR.
Based on this technological framework, the next chapter presents the design,
implementation, and evaluation of the proposed grading application.

\end{document}
